در این قسمت با توجه به خواسته مسئله، از خواص تابع دلتا و فیلتر پایین‌گذر برای یافتن خروجی استفاده می‌کنیم. ابتدا تبدیل فوریه زمان-گسسته (\lr{DTFT}) سیگنال ورودی را می‌نویسیم. می‌دانیم که تبدیل فوریه یک قطار ضربه دوره‌ای، یک قطار ضربه در حوزه فرکانس است:
\begin{flalign*}
    X(e^{j\omega}) &= \frac{2\pi}{16} \sum_{k=-\infty}^{\infty} \delta\left(\omega - \frac{2\pi k}{16}\right) &&
\end{flalign*}

پاسخ فرکانسی سیستم داده شده متشکل از یک فیلتر پایین‌گذر ایده‌آل با فرکانس قطع $\omega_c = \frac{3\pi}{16}$ است که یک تاخیر زمانی به اندازه ۳ نمونه (به دلیل وجود ترم $e^{-j3\omega}$) در آن اعمال شده است.

برای یافتن تبدیل فوریه خروجی، باید $X(e^{j\omega})$ را در پاسخ فرکانسی $H(e^{j\omega})$ ضرب کنیم. طبق خاصیت تابع دلتا در ضرب ($f(x)\delta(x-a) = f(a)\delta(x-a)$)، مقدار پاسخ فرکانسی سیستم تنها در فرکانس‌های مضرب $\frac{2\pi}{16}$ نمونه‌برداری می‌شود:
\begin{flalign*}
    Y(e^{j\omega}) &= X(e^{j\omega}) H(e^{j\omega}) \nonumber \\
    &= \frac{2\pi}{16} \sum_{k=-\infty}^{\infty} H\left(e^{j \frac{2\pi k}{16}}\right) \delta\left(\omega - \frac{2\pi k}{16}\right) &&
\end{flalign*}

با توجه به ویژگی فیلتر پایین‌گذر سیستم، این فیلتر تنها فرکانس‌های درون بازه $|\omega| < \frac{3\pi}{16}$ را از خود عبور می‌دهد و بقیه را صفر می‌کند. بررسی مضارب $\frac{2\pi k}{16}$ نشان می‌دهد که در بازه اصلی فرکانسی، تنها به ازای $k = -1, 0, 1$ فرکانس‌ها درون ناحیه عبور فیلتر قرار می‌گیرند و برای سایر مقادیر $k$، مقدار تابع $H$ صفر خواهد بود.
بنابراین حاصل ضرب به صورت زیر ساده می‌شود:
\begin{flalign*}
    Y(e^{j\omega}) &= \frac{2\pi}{16} \sum_{k=-1}^{1} e^{-j 3 \left(\frac{2\pi k}{16}\right)} \delta\left(\omega - \frac{2\pi k}{16}\right) &&
\end{flalign*}

حال با گرفتن عکس تبدیل فوریه زمان-گسسته (\lr{Inverse DTFT})، سیگنال خروجی در حوزه زمان به دست می‌آید. از آنجا که عکس تبدیل فوریه ضربه $\delta(\omega - \omega_0)$ برابر با $\frac{1}{2\pi}e^{j\omega_0 n}$ است، داریم:
\begin{flalign*}
    y[n] &= \frac{1}{16} \sum_{k=-1}^{1} e^{-j \frac{6\pi k}{16}} e^{j \frac{2\pi k}{16} n} \nonumber \\
    &= \frac{1}{16} \left( e^{-j \frac{2\pi}{16} (n-3)} + 1 + e^{j \frac{2\pi}{16} (n-3)} \right) &&
\end{flalign*}

در نهایت با استفاده از اتحاد اویلر ($e^{j\theta} + e^{-j\theta} = 2\cos(\theta)$)، پاسخ سیستم را به فرم ساده‌تر و نهایی زیر می‌نویسیم:
\begin{flalign*}
    y[n] &= \frac{1}{16} \left( 1 + 2 \cos\left(\frac{\pi}{8} (n-3)\right) \right) &&
\end{flalign*}